%% toc: P1 1 1.3,1.4
\section{Proofs for formula schemas and codes}
\label{sec:annex-p01-w02-p01-004-schemas-codes}

Throughout, \(A\) is the propositional alphabet, \(A^{+}=A\cup\Xi\),
\(\mathsf{Fm}=\mathsf{Fm}_{\mathsf{CPL}}\), and the notation is that of
\cref{sec:p01-w02-p01-004-schemas-codes}. The token map
\(\operatorname{tk}:A^{+}\to A\) sends \(p_n\) to \(p_{2n}\), \(\xi_k\) to \(p_{2k+1}\), and
every other token to itself; \(\operatorname{tr}\) applies \(\operatorname{tk}\) at every
position of a string, so \(\operatorname{tr}(w)\) has the length of \(w\) and
\(\operatorname{tr}(w)(i)=\operatorname{tk}(w(i))\).

\begin{annexitem}{annex:p01-w02-p01-004-schema-readability}{thm:p01-w02-p01-004-schema-readability}{Unique readability for schemas}
\begin{proof}
\emph{Step 1: \(\operatorname{tr}\) is injective.} The map \(\operatorname{tk}\) is
injective: it is the identity on the non-variable tokens, it sends the
remaining tokens to propositional variables, which are distinct from every
non-variable token, and \(p_{2n}\ne p_{2k+1}\) for all \(n,k\), while
\(n\mapsto 2n\) and \(k\mapsto 2k+1\) are injective. If
\(\operatorname{tr}(w)=\operatorname{tr}(w')\), the two strings have equal
length and \(\operatorname{tk}(w(i))=\operatorname{tk}(w'(i))\) at each position, so \(w=w'\).

\emph{Step 2: \(\operatorname{tr}\) commutes with the formation clauses.}
Since \(\operatorname{tr}\) acts position by position, it preserves
concatenation; as \(\operatorname{tk}\) fixes \(\neg\), the binary connectives and both
parentheses, \(\operatorname{tr}(\neg\mathfrak s)=\neg\operatorname{tr}(\mathfrak
s)\) and \(\operatorname{tr}((\mathfrak s\circ\mathfrak t))=
(\operatorname{tr}(\mathfrak s)\circ\operatorname{tr}(\mathfrak t))\).

\emph{Step 3: \(\operatorname{tr}\) maps \(\mathrm{Sch}\) into \(\mathsf{Fm}\).}
Let \(X\) be the set of strings \(w\) over \(A^{+}\) with
\(\operatorname{tr}(w)\in\mathsf{Fm}\). It contains \(\xi_k\) and \(p_n\),
whose images are the variables \(p_{2k+1}\) and \(p_{2n}\), and \(\bottom\).
By Step 2 and the formation clauses of \(\mathsf{Fm}\), \(X\) is closed under
the negation and binary clauses of \cref{def:p01-w02-p01-004-schema}. As
\(\mathrm{Sch}\) is the least set with these closure properties,
\(\mathrm{Sch}\subseteq X\).

\emph{Step 4: \(\operatorname{tr}\) maps \(\mathrm{Sch}\) onto \(\mathsf{Fm}\).}
Let \(Y=\{\operatorname{tr}(\mathfrak s):\mathfrak s\in\mathrm{Sch}\}\). For
every \(m\), \(p_m\in Y\): if \(m=2n\) it is \(\operatorname{tr}(p_n)\), and if
\(m=2k+1\) it is \(\operatorname{tr}(\xi_k)\). Also
\(\bottom=\operatorname{tr}(\bottom)\in Y\). By Step 2, \(Y\) is closed under
\(\neg\) and under the binary clauses. Since \(\mathsf{Fm}\) is the least set
with these properties, \(\mathsf{Fm}\subseteq Y\); with Step 3,
\(\operatorname{tr}\) is a bijection from \(\mathrm{Sch}\) onto \(\mathsf{Fm}\).

\emph{Step 5: exactly one clause, and unique components.} A schema produced
by the first clause is the single token \(\xi_k\); by the second, the single
token \(p_n\) or \(\bottom\); by the third, a string beginning with \(\neg\);
by the fourth, a string beginning with a left parenthesis. These four sets are
pairwise disjoint because the tokens \(\xi_k\), \(p_n\), \(\bottom\), \(\neg\)
and the left parenthesis are pairwise distinct, and a string beginning with
\(\neg\) or a parenthesis is not one of the listed single tokens. If
\(\neg\mathfrak s=\neg\mathfrak s'\), deleting the first token gives
\(\mathfrak s=\mathfrak s'\). Suppose
\((\mathfrak s_1\circ\mathfrak s_2)=(\mathfrak s_1'\circ'\mathfrak s_2')\) with
all four components in \(\mathrm{Sch}\). Applying \(\operatorname{tr}\) and
Step 2 gives
\((\operatorname{tr}\mathfrak s_1\circ\operatorname{tr}\mathfrak s_2)
=(\operatorname{tr}\mathfrak s_1'\circ'\operatorname{tr}\mathfrak s_2')\), an
equation between binary formulas whose displayed components are formulas by
Step 3. By \cref{thm:prop-unique-readability}, \(\circ=\circ'\),
\(\operatorname{tr}\mathfrak s_1=\operatorname{tr}\mathfrak s_1'\) and
\(\operatorname{tr}\mathfrak s_2=\operatorname{tr}\mathfrak s_2'\). Step 1
gives \(\mathfrak s_1=\mathfrak s_1'\) and \(\mathfrak s_2=\mathfrak s_2'\).

\emph{Step 6: the formulas are the letter-free schemas.} The set
\(\mathrm{Sch}\) contains every \(p_n\) and \(\bottom\) and is closed under
\(\neg\) and the binary clauses, so \(\mathsf{Fm}\subseteq\mathrm{Sch}\) by
leastness of \(\mathsf{Fm}\); a formula is a string over \(A\), so it contains
no schematic letter. Conversely, let \(Z\) be the set of schemas
\(\mathfrak s\) such that \(\operatorname{MV}(\mathfrak s)=\varnothing\)
implies \(\mathfrak s\in\mathsf{Fm}\). Each \(\xi_k\) is in \(Z\) vacuously,
since \(\operatorname{MV}(\xi_k)=\{\xi_k\}\); \(p_n\) and \(\bottom\) are
formulas. If \(\mathfrak s\in Z\) and
\(\operatorname{MV}(\neg\mathfrak s)=\operatorname{MV}(\mathfrak
s)=\varnothing\), then \(\mathfrak s\in\mathsf{Fm}\), so
\(\neg\mathfrak s\in\mathsf{Fm}\). If \(\mathfrak s,\mathfrak t\in Z\) and
\(\operatorname{MV}((\mathfrak s\circ\mathfrak t))=\operatorname{MV}(\mathfrak
s)\cup\operatorname{MV}(\mathfrak t)=\varnothing\), both components are
formulas and so is \((\mathfrak s\circ\mathfrak t)\). Hence \(Z\) is closed
under the four clauses and contains \(\mathrm{Sch}\).
\end{proof}
\end{annexitem}

\begin{annexitem}{annex:p01-w02-p01-004-matching-unique}{thm:p01-w02-p01-004-matching-unique}{Uniqueness of matching}
\begin{proof}
Both clauses are proved by induction on \(\mathfrak s\), which is legitimate
because \(\mathrm{Sch}\) is a least set and, by
\cref{thm:p01-w02-p01-004-schema-readability}, each schema has one formation
clause and unique components.

\emph{Clause 1.} If \(\mathfrak s=\xi_k\), then
\(\mathfrak s[\iota]=\iota(\xi_k)=\iota'(\xi_k)=\mathfrak s[\iota']\) because
\(\xi_k\in\operatorname{MV}(\mathfrak s)\). If \(\mathfrak s\) is \(p_n\) or
\(\bottom\), both instances equal \(\mathfrak s\). If
\(\mathfrak s=\neg\mathfrak s_1\), then
\(\operatorname{MV}(\mathfrak s_1)=\operatorname{MV}(\mathfrak s)\), the
induction hypothesis gives \(\mathfrak s_1[\iota]=\mathfrak s_1[\iota']\), and
prefixing \(\neg\) gives the claim. If
\(\mathfrak s=(\mathfrak s_1\circ\mathfrak s_2)\), then \(\iota\) and
\(\iota'\) agree on \(\operatorname{MV}(\mathfrak s_1)\) and on
\(\operatorname{MV}(\mathfrak s_2)\), both subsets of
\(\operatorname{MV}(\mathfrak s)\); the induction hypothesis gives equal
components, hence equal instances.

\emph{Clause 2.} If \(\mathfrak s=\xi_k\), then
\(\iota(\xi_k)=\mathfrak s[\iota]=\mathfrak s[\iota']=\iota'(\xi_k)\), and
\(\operatorname{MV}(\mathfrak s)=\{\xi_k\}\). If \(\mathfrak s\) is \(p_n\) or
\(\bottom\), then \(\operatorname{MV}(\mathfrak s)=\varnothing\) and there is
nothing to prove. If \(\mathfrak s=\neg\mathfrak s_1\), the hypothesis reads
\(\neg(\mathfrak s_1[\iota])=\neg(\mathfrak s_1[\iota'])\); deleting the first
token gives \(\mathfrak s_1[\iota]=\mathfrak s_1[\iota']\), and the induction
hypothesis gives agreement on
\(\operatorname{MV}(\mathfrak s_1)=\operatorname{MV}(\mathfrak s)\). If
\(\mathfrak s=(\mathfrak s_1\circ\mathfrak s_2)\), the hypothesis reads
\((\mathfrak s_1[\iota]\circ\mathfrak s_2[\iota])
=(\mathfrak s_1[\iota']\circ\mathfrak s_2[\iota'])\). All four components are
formulas, so \cref{thm:prop-unique-readability} gives
\(\mathfrak s_1[\iota]=\mathfrak s_1[\iota']\) and
\(\mathfrak s_2[\iota]=\mathfrak s_2[\iota']\). The induction hypothesis gives
agreement on \(\operatorname{MV}(\mathfrak s_1)\) and on
\(\operatorname{MV}(\mathfrak s_2)\), whose union is
\(\operatorname{MV}(\mathfrak s)\).

The final sentence of the theorem is clause 2: two instantiations producing
the same instance agree on every letter of \(\mathfrak s\).
\end{proof}
\end{annexitem}

\begin{annexitem}{annex:p01-w02-p01-004-matching-correct}{thm:p01-w02-p01-004-matching-correct}{Correctness of matching}
\begin{proof}
The procedure is a recursion on \(\mathfrak s\), well defined by
\cref{thm:p01-w02-p01-004-schema-readability}; at each step the decomposition
of \(\varphi\) it uses is unique by \cref{thm:prop-unique-readability}. We
prove clauses 1 and 2 together by induction on \(\mathfrak s\), for all
formulas \(\varphi\) at once.

\emph{Case \(\mathfrak s=\xi_k\).} The output is \(m=\{\xi_k\mapsto\varphi\}\),
never \(\mathsf{fail}\), so clause 2 holds. Its domain is
\(\{\xi_k\}=\operatorname{MV}(\mathfrak s)\), and every \(\iota\) extending
\(m\) has \(\mathfrak s[\iota]=\iota(\xi_k)=\varphi\).

\emph{Case \(\mathfrak s=p_n\) or \(\mathfrak s=\bottom\).} If the output is not
\(\mathsf{fail}\), then \(\varphi=\mathfrak s\) and \(m\) is empty, with domain
\(\operatorname{MV}(\mathfrak s)=\varnothing\), and \(\mathfrak
s[\iota]=\mathfrak s=\varphi\) for every \(\iota\). If
\(\varphi\in\operatorname{Inst}(\mathfrak s)\), then \(\varphi=\mathfrak
s[\iota]=\mathfrak s\) for some \(\iota\), so the output is not
\(\mathsf{fail}\).

\emph{Case \(\mathfrak s=\neg\mathfrak s_1\).} Clause 1: an output other than
\(\mathsf{fail}\) arises only when \(\varphi=\neg\varphi_1\) and
\(m=\operatorname{match}(\mathfrak s_1,\varphi_1)\). By the induction
hypothesis its domain is \(\operatorname{MV}(\mathfrak
s_1)=\operatorname{MV}(\mathfrak s)\), and for \(\iota\) extending \(m\),
\(\mathfrak s_1[\iota]=\varphi_1\), so
\(\mathfrak s[\iota]=\neg\varphi_1=\varphi\). Clause 2: if
\(\varphi=\mathfrak s[\iota]=\neg(\mathfrak s_1[\iota])\), then \(\varphi\) is
a negation with immediate subformula \(\varphi_1=\mathfrak s_1[\iota]\in
\operatorname{Inst}(\mathfrak s_1)\), and the induction hypothesis shows that
\(\operatorname{match}(\mathfrak s_1,\varphi_1)\), which is the output, is
not \(\mathsf{fail}\).

\emph{Case \(\mathfrak s=(\mathfrak s_1\circ\mathfrak s_2)\).} Clause 1: an
output other than \(\mathsf{fail}\) arises only when
\(\varphi=(\varphi_1\circ\varphi_2)\), \(m_1\) and \(m_2\) are not
\(\mathsf{fail}\), they agree on every common letter, and \(m=m_1\cup m_2\),
which is then a function. By the induction hypothesis its domain is
\(\operatorname{MV}(\mathfrak s_1)\cup\operatorname{MV}(\mathfrak
s_2)=\operatorname{MV}(\mathfrak s)\). An \(\iota\) extending \(m\) extends
\(m_1\) and \(m_2\), so \(\mathfrak s_1[\iota]=\varphi_1\) and
\(\mathfrak s_2[\iota]=\varphi_2\), and \(\mathfrak s[\iota]=\varphi\).
Clause 2: let \(\varphi=\mathfrak s[\iota]=(\mathfrak s_1[\iota]\circ
\mathfrak s_2[\iota])\). By \cref{thm:prop-unique-readability} the main
connective of \(\varphi\) is \(\circ\) and its components are
\(\varphi_i=\mathfrak s_i[\iota]\in\operatorname{Inst}(\mathfrak s_i)\). By the
induction hypothesis (clause 2) \(m_i=\operatorname{match}(\mathfrak
s_i,\varphi_i)\ne\mathsf{fail}\). By the induction hypothesis (clause 1) every
instantiation \(\iota_i\) extending \(m_i\) satisfies \(\mathfrak
s_i[\iota_i]=\varphi_i=\mathfrak s_i[\iota]\), so
\cref{thm:p01-w02-p01-004-matching-unique}, clause 2, gives
\(m_i=\iota\mathbin{\upharpoonright}\operatorname{MV}(\mathfrak s_i)\). Hence \(m_1\) and
\(m_2\) both agree with \(\iota\), so they agree on every common letter, and
the output is \(m_1\cup m_2\ne\mathsf{fail}\).

\emph{Consequences.} If the output is \(m\ne\mathsf{fail}\), extend \(m\) to an
instantiation, for instance by sending every letter outside its domain to
\(\bottom\); clause 1 shows \(\varphi\in\operatorname{Inst}(\mathfrak s)\).
The converse is clause 2. If \(\mathfrak s[\iota]=\varphi\) and
\(m=\operatorname{match}(\mathfrak s,\varphi)\), any extension \(\iota'\) of
\(m\) satisfies \(\mathfrak s[\iota']=\varphi=\mathfrak s[\iota]\), so
\(\iota\) agrees with \(\iota'\), hence with \(m\), on
\(\operatorname{MV}(\mathfrak s)\), the domain of \(m\), by
\cref{thm:p01-w02-p01-004-matching-unique}.
\end{proof}
\end{annexitem}

\begin{annexitem}{annex:p01-w02-p01-004-instance-count}{thm:p01-w02-p01-004-instance-count}{Size of instance sets}
\begin{proof}
\emph{Clause 1.} By Step 6 of the proof of
\cref{thm:p01-w02-p01-004-schema-readability}, \(\mathfrak s\in\mathsf{Fm}\).
We show by induction on \(\mathfrak s\in\mathrm{Sch}\) that if
\(\operatorname{MV}(\mathfrak s)=\varnothing\), then
\(\mathfrak s[\iota]=\mathfrak s\) for every instantiation \(\iota\). Every
letter occurring in an immediate subschema of \(\mathfrak s\) occurs in
\(\mathfrak s\), so \(\operatorname{MV}(\mathfrak s_1)\subseteq
\operatorname{MV}(\mathfrak s)\) when \(\mathfrak s=\neg\mathfrak s_1\), and
\(\operatorname{MV}(\mathfrak s_1)\cup\operatorname{MV}(\mathfrak s_2)\subseteq
\operatorname{MV}(\mathfrak s)\) when \(\mathfrak s=(\mathfrak s_1\circ
\mathfrak s_2)\).
\emph{Case \(\mathfrak s=\xi_k\).} Then \(\operatorname{MV}(\mathfrak
s)=\{\xi_k\}\ne\varnothing\), so the hypothesis fails and there is nothing to
prove.
\emph{Case \(\mathfrak s=p_n\).} By the instantiation clause for \(p_n\),
\(\mathfrak s[\iota]=p_n=\mathfrak s\).
\emph{Case \(\mathfrak s=\bottom\).} By the clause for \(\bottom\),
\(\mathfrak s[\iota]=\bottom=\mathfrak s\).
\emph{Case \(\mathfrak s=\neg\mathfrak s_1\).} Here
\(\operatorname{MV}(\mathfrak s_1)\subseteq\operatorname{MV}(\mathfrak
s)=\varnothing\), so the induction hypothesis gives \(\mathfrak
s_1[\iota]=\mathfrak s_1\), and the negation clause gives
\(\mathfrak s[\iota]=\neg(\mathfrak s_1[\iota])=\neg\mathfrak s_1=\mathfrak s\).
\emph{Case \(\mathfrak s=(\mathfrak s_1\circ\mathfrak s_2)\).} Here
\(\operatorname{MV}(\mathfrak s_1)\) and \(\operatorname{MV}(\mathfrak s_2)\)
are subsets of \(\operatorname{MV}(\mathfrak s)=\varnothing\), so the
induction hypothesis gives \(\mathfrak s_1[\iota]=\mathfrak s_1\) and
\(\mathfrak s_2[\iota]=\mathfrak s_2\), and the binary clause gives
\(\mathfrak s[\iota]=(\mathfrak s_1[\iota]\circ\mathfrak s_2[\iota])
=(\mathfrak s_1\circ\mathfrak s_2)=\mathfrak s\).
These cases exhaust the formation clauses of
\cref{def:p01-w02-p01-004-schema}, so every instance of \(\mathfrak s\) is
\(\mathfrak s\). Instantiations exist (the map sending every letter to
\(\bottom\) is one), so \(\operatorname{Inst}(\mathfrak s)=\{\mathfrak s\}\).

\emph{Clause 2.} Fix \(\xi_k\in\operatorname{MV}(\mathfrak s)\). For each
\(n\in\N\) let \(\iota_n\) send every letter to \(p_n\). If
\(\mathfrak s[\iota_n]=\mathfrak s[\iota_{n'}]\), then
\cref{thm:p01-w02-p01-004-matching-unique}, clause 2, gives
\(p_n=\iota_n(\xi_k)=\iota_{n'}(\xi_k)=p_{n'}\), so \(n=n'\). Hence the
instances \(\mathfrak s[\iota_n]\) are pairwise distinct and
\(\operatorname{Inst}(\mathfrak s)\) is infinite. By
\cref{thm:p01-w02-p01-004-code-properties}, clause 3, the code \(\#\) maps
\(\operatorname{Inst}(\mathfrak s)\subseteq\mathsf{Fm}\) injectively onto an
infinite subset \(I\) of \(\N\). Enumerating \(I\) in increasing order (the
\(j\)-th member is the least member of \(I\) exceeding the earlier ones, which
exists because \(I\) is infinite) gives a bijection \(\N\to I\). Composing with
the inverse of \(\#\) on \(I\) gives a bijection
\(\N\to\operatorname{Inst}(\mathfrak s)\).
\end{proof}
\end{annexitem}

\begin{annexitem}{annex:p01-w02-p01-004-substitution-link}{thm:p01-w02-p01-004-substitution-link}{Instances as uniform-substitution instances}
\begin{proof}
\emph{Clause 1.} \(\operatorname{tr}(\mathfrak s)\in\mathsf{Fm}\) is Step 3 of
the proof of \cref{thm:p01-w02-p01-004-schema-readability}. The identity is
proved by induction on \(\mathfrak s\), using the recursion of
\cref{def:p02-uniform-substitution} on the right-hand side. For
\(\mathfrak s=\xi_k\):
\(\operatorname{tr}(\xi_k)[\upsilon_\iota]=p_{2k+1}[\upsilon_\iota]
=\upsilon_\iota(p_{2k+1})=\iota(\xi_k)=\xi_k[\iota]\). For
\(\mathfrak s=p_n\):
\(p_{2n}[\upsilon_\iota]=p_n=p_n[\iota]\). For \(\bottom\) both sides are
\(\bottom\). For \(\mathfrak s=\neg\mathfrak s_1\), Step 2 of that proof gives
\(\operatorname{tr}(\mathfrak s)[\upsilon_\iota]
=(\neg\operatorname{tr}(\mathfrak s_1))[\upsilon_\iota]
=\neg(\operatorname{tr}(\mathfrak s_1)[\upsilon_\iota])
=\neg(\mathfrak s_1[\iota])=\mathfrak s[\iota]\) by the induction hypothesis.
For \(\mathfrak s=(\mathfrak s_1\circ\mathfrak s_2)\), Step 2 of that
proof, the binary clause of \cref{def:p02-uniform-substitution}, the induction
hypothesis for \(\mathfrak s_1\) and \(\mathfrak s_2\), and the binary
instantiation clause give
\(\operatorname{tr}(\mathfrak s)[\upsilon_\iota]
=(\operatorname{tr}(\mathfrak s_1)\circ\operatorname{tr}(\mathfrak
s_2))[\upsilon_\iota]
=(\operatorname{tr}(\mathfrak s_1)[\upsilon_\iota]\circ
\operatorname{tr}(\mathfrak s_2)[\upsilon_\iota])
=(\mathfrak s_1[\iota]\circ\mathfrak s_2[\iota])
=\mathfrak s[\iota]\).

\emph{Clause 2.} Each \(\upsilon_\iota\) satisfies
\(\upsilon_\iota(p_{2n})=p_n\), so clause 1 gives the inclusion from left to
right. Conversely, let \(\upsilon\) be a uniform substitution with
\(\upsilon(p_{2n})=p_n\) for all \(n\), and define \(\iota(\xi_k)=
\upsilon(p_{2k+1})\). Then \(\upsilon\) and \(\upsilon_\iota\) agree on every
propositional variable, so they are the same function and
\(\operatorname{tr}(\mathfrak s)[\upsilon]=\mathfrak s[\iota]\in
\operatorname{Inst}(\mathfrak s)\).

\emph{Clause 3.} Let \(\varphi\in\operatorname{Inst}(\mathfrak s)\), say
\(\varphi=\mathfrak s[\iota]\). By clause 1,
\(\varphi=\operatorname{tr}(\mathfrak s)[\upsilon_\iota]\). By
\cref{thm:p02-uniform-substitution}, validity of
\(\operatorname{tr}(\mathfrak s)\) implies validity of
\(\operatorname{tr}(\mathfrak s)[\upsilon_\iota]\).
\end{proof}
\end{annexitem}

\begin{annexitem}{annex:p01-w02-p01-004-code-properties}{thm:p01-w02-p01-004-code-properties}{Properties of the code}
\begin{proof}
\emph{Clause 1.} Put \(t_s=s(s+1)/2\), so \(\pi(a,b)=t_{a+b}+b\). For
\(s\ge1\), \((s+1)/2\ge1\) gives \(t_s\ge s\), and \(t_0=0\). Hence
\(\pi(a,b)\ge(a+b)+b\), which is at least \(a\) and at least \(b\).

\emph{Clause 2.} Every line of \cref{def:p01-w02-p01-004-code} has the form
\(1+\pi(\cdot,\cdot)\) with a value in \(\N\).

\emph{Clause 3.} We prove by induction on \(\varphi\) that for every formula
\(\psi\), \(\#\varphi=\#\psi\) implies \(\varphi=\psi\). Write
\(\#\varphi=1+\pi(k,x)\) and \(\#\psi=1+\pi(k',x')\) according to the unique
lines that apply (\cref{thm:prop-unique-readability}). From
\(\#\varphi=\#\psi\), injectivity of \(\pi\) (\cref{sec:strings-01}) gives
\(k=k'\) and \(x=x'\). The tag determines the formation clause and, for a
binary formula, the connective; so \(\varphi\) and \(\psi\) are formed by the
same clause. Tag \(0\): \(\varphi=p_x=\psi\). Tag \(1\): both are \(\bottom\).
Tag \(2\): \(\varphi=\neg\varphi_1\), \(\psi=\neg\psi_1\) and
\(\#\varphi_1=x=\#\psi_1\); the induction hypothesis for \(\varphi_1\) gives
\(\varphi_1=\psi_1\). Tags \(3,4,5\): \(\varphi=(\varphi_1\circ\varphi_2)\)
and \(\psi=(\psi_1\circ\psi_2)\) with the same \(\circ\), and
\(\pi(\#\varphi_1,\#\varphi_2)=x=\pi(\#\psi_1,\#\psi_2)\); injectivity of
\(\pi\) gives \(\#\varphi_i=\#\psi_i\), and the induction hypothesis for
\(\varphi_1\) and \(\varphi_2\) gives \(\varphi_i=\psi_i\).

\emph{Clause 4.} For \(\neg\varphi\), clause 1 gives
\(\#\varphi\le\pi(2,\#\varphi)<1+\pi(2,\#\varphi)=\#\neg\varphi\). For a
binary formula with tag \(k\), clause 1 applied twice gives
\(\#\varphi_i\le\pi(\#\varphi_1,\#\varphi_2)\le\pi(k,\pi(\#\varphi_1,\#\varphi_2))
<\#(\varphi_1\circ\varphi_2)\) for \(i=1,2\).
\end{proof}
\end{annexitem}

\begin{annexitem}{annex:p01-w02-p01-004-decoding-correct}{thm:p01-w02-p01-004-decoding-correct}{Correctness of decoding}
\begin{proof}
The inverse \(\pi^{-1}\) used by the procedure is correct by the bijectivity
argument of \cref{sec:strings-01}, and the procedure terminates on every input
by the bound stated with it. Termination makes
\(\operatorname{dec}\) a total function on \(\N\) with values in
\(\mathsf{Fm}\cup\{\mathsf{none}\}\).

\emph{First claim, by induction on \(\varphi\).} Write
\(\#\varphi-1=\pi(k,x)\) according to \cref{def:p01-w02-p01-004-code};
injectivity of \(\pi\) gives \(\pi^{-1}(\#\varphi-1)=(k,x)\), and
\(\#\varphi\ne0\). If \(\varphi=p_n\), then \(k=0\), \(x=n\), and the output is
\(p_n\). If \(\varphi=\bottom\), then \(k=1\), \(x=0\), and the output is
\(\bottom\). If \(\varphi=\neg\varphi_1\), then \(k=2\), \(x=\#\varphi_1\), and
by the induction hypothesis \(\operatorname{dec}(x)=\varphi_1\), so the
output is \(\neg\varphi_1\). If \(\varphi=(\varphi_1\circ\varphi_2)\), then
\(k\) is the tag of \(\circ\), \(x=\pi(\#\varphi_1,\#\varphi_2)\),
\(\pi^{-1}(x)=(\#\varphi_1,\#\varphi_2)\), and by the induction hypothesis the
output is \((\varphi_1\circ_k\varphi_2)=\varphi\).

\emph{Second claim, by strong induction on \(c\).} Suppose
\(\operatorname{dec}(c)=\varphi\) is a formula. Then \(c\ne0\); let
\((k,x)=\pi^{-1}(c-1)\), so \(c=1+\pi(k,x)\), and \(k\le5\) because tags
\(k\ge6\) return \(\mathsf{none}\). If \(k=0\), \(\varphi=p_x\) and
\(\#p_x=1+\pi(0,x)=c\). If \(k=1\), the output is a formula only when \(x=0\),
and then \(\varphi=\bottom\) with \(\#\bottom=1+\pi(1,0)=c\). If \(k=2\),
\(\operatorname{dec}(x)=\varphi_1\) is a formula and \(\varphi=\neg\varphi_1\);
since \(x<c\), the induction hypothesis gives \(\#\varphi_1=x\), so
\(\#\varphi=1+\pi(2,x)=c\). If \(k\in\{3,4,5\}\), with \((a,b)=\pi^{-1}(x)\),
the outputs \(\operatorname{dec}(a)=\varphi_1\) and
\(\operatorname{dec}(b)=\varphi_2\) are formulas and
\(\varphi=(\varphi_1\circ_k\varphi_2)\); as \(a,b\le x<c\), the induction
hypothesis gives \(\#\varphi_1=a\) and \(\#\varphi_2=b\), so
\(\#\varphi=1+\pi(k,\pi(a,b))=1+\pi(k,x)=c\).

\emph{Consequences.} If \(c=\#\varphi\), the first claim gives
\(\operatorname{dec}(c)=\varphi\ne\mathsf{none}\). If
\(\operatorname{dec}(c)\ne\mathsf{none}\), it is a formula \(\varphi\), and the
second claim gives \(c=\#\varphi\). The two claims say that
\(\operatorname{dec}\circ\#\) is the identity on \(\mathsf{Fm}\) and
\(\#\circ\operatorname{dec}\) is the identity on the set of codes.
\end{proof}
\end{annexitem}

\begin{annexitem}{annex:p01-w02-p01-004-distinction}{thm:p01-w02-p01-004-distinction}{Schemas stand for families; codes name single formulas}
\begin{proof}
\emph{Clause 1.} By \cref{thm:p01-w02-p01-004-decoding-correct},
\(\operatorname{dec}(\#\varphi)=\varphi\), and any \(c\) with
\(\operatorname{dec}(c)=\varphi\) equals \(\#\varphi\). Since
\(\operatorname{dec}\) is a function, a number \(c\) with
\(\operatorname{dec}(c)\ne\mathsf{none}\) is correlated with the one formula
\(\operatorname{dec}(c)\).

\emph{Clause 2.} If \(\mathfrak s\) is a formula, it is a string over \(A\),
so \(\operatorname{MV}(\mathfrak s)=\varnothing\), and
\cref{thm:p01-w02-p01-004-instance-count}, clause 1, gives
\(\operatorname{Inst}(\mathfrak s)=\{\mathfrak s\}\). Conversely, if
\(\operatorname{Inst}(\mathfrak s)\) has exactly one member, it is finite, so
by clause 2 of that theorem \(\operatorname{MV}(\mathfrak s)=\varnothing\),
and by its clause 1 \(\mathfrak s\) is a formula.

\emph{Clause 3.} For a map \(\lambda:\Xi\to\Xi\), let \(\mathfrak s^{\lambda}\)
be the string obtained from \(\mathfrak s\) by replacing each letter
occurrence \(\xi\) by \(\lambda(\xi)\). Induction on \(\mathfrak s\) shows that
\(\mathfrak s^{\lambda}\in\mathrm{Sch}\) and
\(\mathfrak s^{\lambda}[\iota]=\mathfrak s[\iota\circ\lambda]\) for every
instantiation \(\iota\): for \(\xi_k\) both sides are
\(\iota(\lambda(\xi_k))\); for \(p_n\) and \(\bottom\) both sides are the
token; and the renaming commutes with \(\neg\) and with the binary clauses.

Now let \(\xi_k\in\operatorname{MV}(\mathfrak s)\). Because
\(\operatorname{MV}(\mathfrak s)\) is finite and \(\Xi\) is infinite, choose
\(\xi_j\notin\operatorname{MV}(\mathfrak s)\), and let \(\lambda\) exchange
\(\xi_k\) and \(\xi_j\) and fix every other letter. Put
\(\mathfrak s'=\mathfrak s^{\lambda}\). At a position where \(\mathfrak s\) has
\(\xi_k\), \(\mathfrak s'\) has \(\xi_j\), so \(\mathfrak s'\ne\mathfrak s\).
Since \(\lambda\circ\lambda\) is the identity, \(\iota\mapsto\iota\circ\lambda\)
is a bijection of the set of instantiations onto itself. Therefore
\(\operatorname{Inst}(\mathfrak s')=\{\mathfrak s[\iota\circ\lambda]:\iota\}
=\{\mathfrak s[\iota']:\iota'\}=\operatorname{Inst}(\mathfrak s)\). The
contrast in the last sentence of the theorem is clause 1 for codes and the
present clause for schemas.
\end{proof}
\end{annexitem}

\begin{annexitem}{annex:p01-w02-p01-004-instance-codes}{thm:p01-w02-p01-004-instance-codes}{The codes of an instance family}
\begin{proof}
If \(c=\#\varphi\) with \(\varphi\in\operatorname{Inst}(\mathfrak s)\), then
\(\operatorname{dec}(c)=\varphi\ne\mathsf{none}\) by
\cref{thm:p01-w02-p01-004-decoding-correct}, and
\(\operatorname{match}(\mathfrak s,\varphi)\ne\mathsf{fail}\) by
\cref{thm:p01-w02-p01-004-matching-correct}, clause 2. Conversely, if
\(\operatorname{dec}(c)=\varphi\) is a formula, then \(c=\#\varphi\) by
\cref{thm:p01-w02-p01-004-decoding-correct}, and if moreover
\(\operatorname{match}(\mathfrak s,\varphi)\ne\mathsf{fail}\), then
\(\varphi\in\operatorname{Inst}(\mathfrak s)\) by
\cref{thm:p01-w02-p01-004-matching-correct}. Both procedures terminate on
every admissible input, by the bounds stated with them, so running the first
on \(c\) and, if it returns a formula, the second on \(\mathfrak s\) and that
formula, answers the membership question for every \(c\).
\end{proof}
\end{annexitem}
